%\section{Introduction}
Many problems in economics and finance involve decision-making that cannot be taken continuously in time, but rather at discrete and irregular intervals. For example, a central bank does not continuously intervene in the foreign exchange market, a (retail) investor facing transaction costs does not rebalance her portfolio at every instant, and a firm does not continuously inject capital into a project for every small fluctuation in profitability. In all these cases, the controller, whether a central bank, investor, or firm, typically waits until an intervention is sufficiently valuable. The controller must therefore decide both \textit{when} to intervene and \textit{what} intervention to apply, resulting in a stochastic impulse control problem.

Between interventions, the state process is assumed to evolve according to a stochastic differential equation (SDE). In this thesis, the uncontrolled state is modelled as a jump-diffusion of the form
\begin{equation*}
    \diff X(t) = \mu(t,X(t^-)) \diff t + \sigma(t,X(t^-)) \diff W(t) + \int_{\R^\ell} \gamma(t,X(t^-),z) \tilde N(\diff t, \diff z),
\end{equation*}
where the Brownian motion $W$ models continuous uncertainty and the compensated Poisson random measure $\tilde N$ models jumps. 

An impulse control, modeled through a double-sequence $\alpha=\{(\tau_j,\zeta_j)\}_{j \geq 1}$, controls the jump diffusion process at intervention times $\tau_j$ through an impulse $\zeta_j$. If the performance of an impulse control is quantified with the performance criterion $J$, the impulse control problem translates into the maximization problem
\begin{equation*}
    v(t,x) = \sup_{\alpha \in \mathcal{V}}J^{(\alpha)}(t,x).
\end{equation*} 

There exist two main approaches to impulse control problems. The first is the dynamic programming approach, based on Bellman's principle of optimality \cite{bellman1957dynamic}. In this formulation, the value function is characterised recursively. That is, after the next decision time, the problem restarts from the new state. For impulse control, this leads to an approximation by a sequence of optimal stopping problems, as described by \textcite{oksendal2007applied}. The second way to solve the impulse-control problem is through the parabolic Hamilton-Jacobi-Bellman Quasi-Variational Inequality (HJBQVI), or simply QVI since continuous controls are not considered in this thesis. The QVI approach was first introduced by \textcite{bensoussan1975nouvelles}, where the authors characterise the value function through the free-boundary problem
\begin{equation}
    \min\{ \underbrace{- (\partial_t v + \InfinitesimalGenerator v + f)}_{\text{No intervention}}, \underbrace{v - \InterveneOper v}_{\text{Intervene now}}\} = 0, 
    \tag{QVI}
    \label{eqn: QVI intro}
\end{equation}
Here, $\InfinitesimalGenerator$ is the infinitesimal generator of the uncontrolled state process, and $\InterveneOper$ is the nonlocal intervention operator. The first term in \eqref{eqn: QVI intro} corresponds to continuation without intervention, while the second term corresponds to immediate intervention. 

In general, the value function $v$ is not expected to be smooth as fixed intervention costs may create (non-differentiable) kinks at the intervention boundary. Moreover, as the underlying stochastic process is a jump-diffusion, the generator $\InfinitesimalGenerator$ contains a non-local and, possibly, singular integral, which further increases the complexity of the problem. As a consequence, classical solutions are either too restrictive or may not even exist in the specified domain, and one instead opts for a weaker solution in the form of \textit{viscosity solutions}. The theory of viscosity solutions was first developed by \textcite{crandall1983viscosity} for local first-order partial differential equations (PDEs) and later extended to local second-order PDEs by \textcite{crandall1992user}. The nonlocal viscosity theory for partial integro-differential equations (PIDEs) is treated in \textcite{barles2008second}, as well as \textcite{jakobsen2006maximum}. For impulse control of jump-diffusions, \textcite{seydel2010generalexistenceuniquenessviscosity} gives the central existence and uniqueness result for the elliptic and parabolic HJBQVIs used in this thesis. 

Closed-form solutions are rarely available for impulse control problems, so numerical methods are typically required. However, \eqref{eqn: QVI intro} is a nonlinear equation with a nonlocal intervention operator, making numerical approximation difficult. While classical convergence arguments of \textcite{barlessouganidis1991convergence} prove that monotone, stable, and consistent schemes converge to the correct solution, the presence of a nonlocal intervention operator requires additional analysis. An extension of the Barles--Souganidis theorem to HJBQVIs is provided in \cite{azimzadeh2018impulsecontrolfinancenumerical,azimzadeh2018convergence,azimzadeh2016weakly}. Nevertheless, finite-difference schemes scale poorly with dimension, which limits their applicability to high-dimensional portfolio problems.

This motivates the second computational direction covered in this thesis. Deep reinforcement learning (DRL) methods have been widely applied in finance, including portfolio optimisation by \textcite{JXL17}, derivative hedging by \textcite{buehler2019deep, buehler2022deep}, as well as trading applications by \textcite{jain2025impulse} and \textcite{anchuri2026rammstein}. In this thesis, two DRL approaches are used to approximate impulse control problems. The value-based approach builds on Double Deep Q-Networks \cite{hasselt2010double, van2016deep}, while the policy-based approach uses Proximal Policy Optimisation \cite{schulman2017proximal} combined with self-imitation learning \cite{oh2018self}.

This thesis focuses on two particular applications of impulse controls in finance: Optimal foreign exchange (FEX) intervention and dynamic portfolio allocation with transaction costs. In the FEX intervention problem, a central bank uses costly market interventions to keep an exchange rate close to a target level. This problem is one of the oldest examples of impulse control in finance, first studied in \cite{mundaca1998optimal, cadenillas1999optimal, cadenillas2000classical}. In the portfolio allocation problem, continuous rebalancing is no longer optimal when trading is costly. Instead, the portfolio manager waits until the portfolio has moved sufficiently far away from the desired allocation before rebalancing. This problem is studied in \cite{EH88, MP95, K98}. 

\section*{Thesis contribution}
This thesis studies known theory on finite-horizon impulse control problems for jump-diffusion processes in a risk-neutral objective setting, in the sense that the controller maximises an expected performance criterion rather than an expected utility. The theoretical part formulates the impulse-control problem, relates it to the dynamic programming principle, and studies the corresponding QVI using viscosity solutions. Existence and uniqueness results are mainly based on \textcite{seydel2010generalexistenceuniquenessviscosity}. However, while Seydel states the parabolic comparison theorem, the proof is concise and largely refers to the elliptic argument. This thesis provides a detailed proof of the comparison result for parabolic QVIs in the finite-horizon and jump-diffusion settings.

In addition, Appendix \ref{appendix: Auxiliary Results for Dynamic Programming} derives a local dynamic-programming approximation in the continuation region. Under suitable smoothness and growth conditions, Proposition \ref{prop: small-h approximation} establishes a residual bound of order $O(h^{3/2})$. To the best of our knowledge, this specific result is not explicitly proven in the references used in this thesis. This result is then used to bridge the impulse-control QVI to the Bellman equations, the latter of which is approximated using value-based DRL methods.

From a numerical perspective, this thesis provides an empirical study of how value-based and policy-based DRL algorithms approximate finite-horizon impulse control problems and intervention regions. In particular, the DRL models are compared against a finite-difference scheme in the foreign exchange (FEX) intervention problem. In the second numerical experiment, the Proximal Policy Optimisation algorithm is applied to a dynamic portfolio allocation problem with transaction costs, where trading decisions are modelled as impulse controls. As such, this thesis aims to provide retail investors with a practical framework to support informed portfolio rebalancing decisions.

\section*{Contribution to sustainability}
Although the main focus of this thesis is the application of impulse control in finance, the theory is applicable in several other fields, including optimal harvesting in forest management \cite{ieda2015implicit}, river restoration \cite{YOSHIOKA20192182}, animal population management \cite{yaegashi2019exact}, and life insurance \cite{azimzadeh2018impulsecontrolfinancenumerical, lu2024semi}. For instance in forest management, impulse control can be used to model optimal harvesting, replanting, or fire-prevention interventions, while applications in life insurance can improve retirement plans. These examples illustrate how impulse controls can help in cases involving responsible resource use, ecosystem protection, and long-term economic planning. As such, they are relevant to several Sustainable Development Goals, including Decent Work and Economic Growth (SDG 8), Industry, Innovation and Infrastructure (SDG 9), Responsible Consumption and Production (SDG 12), Climate Action (SDG 13), and Life on Land (SDG 15). More generally, impulse-control models can support transparent and quantitatively disciplined decision-making in systems where interventions are costly, uncertain, and socially consequential.

\section*{Thesis structure}
The thesis is structured in the following way.
\begin{itemize}
    \item Chapter \ref{section: Preliminaries} introduces the stochastic preliminaries used throughout the thesis, such as filtered probability space, jump-diffusion SDEs, infinitesimal generators, and stopping times.
    \item Chapter \ref{section: Impulse Control of Jump-Diffusion Processes} formulates the impulse-control problem for jump-diffusion processes together with a verification theorem and the dynamic programming principle. 
    \item Chapter \ref{section: Viscosity Solution for Impulse Control: Existence} studies the existence of viscosity solutions for QVIs. The chapter introduces additional assumptions and proves that the value function is a viscosity solution. 
    \item Chapter \ref{section: Viscosity Solution for Impulse Control: Uniqueness} proves the uniqueness of the viscosity solution by means of the comparison-principle for the parabolic QVI. 
    \item Chapter \ref{section: Numerical solution of the QVI} presents a finite-difference scheme used to approximate the QVI in low-dimensional problems. 
    \item Chapter \ref{section: Deep Reinforcement Learning for Impulse Control Problems} gives a brief overview of reinforcement learning and applies DRL two algorithms on the impulse control problem.
    \item Chapter \ref{section: Numerical Experiment I: Impulse Control of the Foreign Exchange Rate} approximates the FEX rate problem using a finite-difference scheme and compares the scheme against DRL algorithms. As such, this chapter can be viewed as numerical benchmark for the DRL methods. 
    \item Chapter \ref{section: Multi-dimensional Dynamic Portfolio Allocation} studies the use of impulse controls in dynamic portfolio allocation and uses the Proximal Policy Optimisation algorithm to approximate a trading policy in a portfolio consisting of 20 risky assets. 
    \item Chapter \ref{section: concluding remarks} discusses the theoretical and numerical findings, outlines possible directions for future work, and concludes the thesis. 
\end{itemize}

% \subsection*{Relevant courses}
% Although this part is not strictly required, we provide a list of courses we found useful for this thesis. Courses taken at the Norwegian University of Science and Technology (NTNU) and Ecole Polytechnique Fédérale de Lausanne (EPFL) are presented separately.

% Useful courses taken at NTNU:
% \begin{itemize}
%     \item \textit{TMA4220 – Numerical Solution of Partial Differential Equations Using Element Methods}: Provided understanding of finite-difference schemes. Used for solving QVIs with finite-difference schemes.
%     \item \textit{TMA4215 – Numerical Mathematics}: Provided understanding of the Fast Fourier Transform algorithm used in the calibration procedure.
%     \item \textit{FIN3009 – Asset Pricing and Portfolio Management}: Provided basic understanding of portfolio management. 
%     \item \textit{TMA4500 – Specialisation project}: Provided understanding for the use of the Fast Fourier Transform in the context of European option pricing. 
% \end{itemize}

% Useful courses taken at EPFL:
% \begin{itemize}
%     \item \textit{FIN417 – Quantitative Risk Management}: Provided basic understanding of risk management and risk metrics for portfolio management. Used in the dynamic portfolio allocation example. 
%     \item \textit{MATH450 – Numerical Integration of Stochastic Differential Equations}: Provided thorough understanding of stochastic calculus and numerical schemes for SDEs. 
%     \item \textit{MATH470 – Martingales in Financial Mathematics}: Provided thorough understanding of stochastic analysis, portfolio hedging, Lévy processes, and hedging using reinforcement learning (Deep Hedging).
%     \item \textit{MATH524 – Nonparameteric estimation and inference}: Provided rigorous results and proofs for universal approximation theorems and statistical consistency proofs for neural networks.
% \end{itemize}